% Generated draft: requires mathematical and source review.
\documentclass[11pt]{article}
\usepackage[margin=25mm]{geometry}
\usepackage{fontspec}
\setmainfont{Latin Modern Roman}
\usepackage{amsmath,amssymb,mathtools,unicode-math}
\setmathfont{Latin Modern Math}
\usepackage{enumitem,needspace,xurl,hyperref}
\setlength{\parindent}{0pt}
\setlength{\parskip}{5pt}
\setlength{\emergencystretch}{3em}
\newcommand{\N}{\mathbb{N}}
\newcommand{\Z}{\mathbb{Z}}
\newcommand{\Q}{\mathbb{Q}}
\newcommand{\R}{\mathbb{R}}
\newcommand{\C}{\mathbb{C}}
\newcommand{\F}{\mathbb{F}}
\newcommand{\A}{\mathbb{A}}
\newcommand{\PP}{\mathbb{P}}
\newcommand{\E}{\mathbb{E}}
\newcommand{\eps}{\varepsilon}
\newcommand{\dd}{\,\mathrm{d}}
\begin{document}
\section*{UnsolvedMath ID: 30005110}
\textbf{Source problem code:} catalogue entry
\subsection{Definitions and mathematical statement}
Let \(k\geq 2\) be an integer. The supplied record uses the terms \(k\)-PDG, \(\widetilde T_k\), and the quantities \(\alpha\) and \(\beta\), but does not define them. In particular, it does not specify the permitted partially directed hyperedges, edge multiplicities, the normalization of \(\alpha\) and \(\beta\), the forbidden configuration \(\widetilde T_k\), or the precise meaning of subgraph containment for these objects. These definitions must be recovered from the cited source before the conjecture has a self-contained mathematical formulation.

Subject to those unresolved source definitions, the target is the following statement, identified in the report excerpt as Conjecture 4: for every integer \(k\geq 2\), there exists an integer \(n_0(k)\geq 1\) such that, for every integer \(n\geq n_0(k)\), every \(n\)-vertex \(k\)-PDG containing no subgraph isomorphic to \(\widetilde T_k\) satisfies
\[
\alpha+(\log_2 3)\beta\leq 1.
\]
Here \(\alpha\) and \(\beta\) are the unspecified quantities associated with that \(k\)-PDG in the source. The dependence of the threshold on \(k\) and the exact inequality are retained; no asymptotic error term is introduced.

The excerpt additionally reports, as Theorem 5, that for each fixed \(k\geq 2\), Conjecture 4 implies its Conjecture 1. It also reports an essential equivalence with Conjecture 1 after replacing the single forbidden configuration by a special finite family \(F_k\). Neither Conjecture 1, the family \(F_k\), nor the meaning of that essential equivalence is supplied, so these contextual claims cannot be formulated precisely here.
\subsection{Short English statement}
For each fixed integer \(k\geq 2\), the conjecture asks whether every sufficiently large \(k\)-PDG avoiding \(\widetilde T_k\) satisfies \(\alpha+(\log_2 3)\beta\leq 1\). The supplied record omits the definitions of the objects, forbidden configuration, and quantities in this inequality, so this is a conditional transcription requiring source clarification and human mathematical review.
\subsection{Sources}
\begin{itemize}
\item Supplied local UnsolvedMath registry record, permanent ID \(30005110\), source problem code \texttt{catalogue entry}, titled ``Extremal Conjecture for Partially Directed Uniform Hypergraphs,'' in the Oberwolfach Reports Open Problems collection. The record gives proposed year 2022 and no named proposer.
\item Supplied report citation: \emph{Combinatorics, Probability and Computing} (2023), Oberwolfach Reports, pp.~1165--1237, DOI \texttt{10.4171/owr/2022/22}. \url{https://doi.org/10.4171/owr/2022/22}. The supplied excerpt identifies the target as Conjecture 4 and its implication to Conjecture 1 as Theorem 5. Oberwolfach Reports are described in the record as documenting mathematical workshops and their open-problem discussions. The supplied citation year 2023 and proposed year 2022 are retained without reconciliation.
\item Supplied literature reference: Dong, Mani, and Zhao, \emph{Enumerating k-SAT functions}. \url{https://arxiv.org/abs/2107.09233}. The report excerpt attributes its \(k=4\) proof to an application of a result of F\"uredi and Melaki, labeled [5] there, on minimum triangular edge density at fixed graph edge density. It attributes the \(k=5\) brute-force verification to an appendix by Nitya Mani and Edward Yu in arXiv-v2 of its reference [4]. The excerpt directs readers to [4] for the precise finite-forbidden-family formulation and its construction. Complete bibliography entries for those numbered references are not supplied; their identification is not assumed here.
\item Literature assessment dated exactly \texttt{2026-08-22}: the supplied curation labels the problem partially solved, reports the conjecture established for \(k=4\) and computationally checked for \(k=5\), and states that no proof for every fixed \(k\) is known. This is supplied dated curation, not an independently verified literature assessment. The record explicitly asks that time-sensitive status and cited claims be verified independently.
\item The record labels its statement \texttt{corrected\_verified} and says that the official report was checked directly and that the clean statement restores definitions and removes extraction artifacts. These are supplied verification claims; no source visit or independent verification is asserted in this draft. The supplied text itself does not include the necessary definitions.
\end{itemize}
\end{document}
